\documentclass[aspectratio=169]{beamer}
\usetheme{default}
\usecolortheme{dove}
\usefonttheme{serif}
\setbeamertemplate{navigation symbols}{}
\setbeamertemplate{footline}{}
\beamertemplatenavigationsymbolsempty

\usepackage{amsmath,amssymb}
\usepackage{graphicx}
\usepackage[most]{tcolorbox}
\usepackage{booktabs}

\graphicspath{{figs/}}

\definecolor{key}{rgb}{0.10,0.35,0.65}
\setbeamercolor{frametitle}{fg=key}

\newcommand{\kb}[1]{\begin{tcolorbox}[colback=key!5,colframe=key!60,boxsep=2pt,left=4pt,right=4pt,top=3pt,bottom=3pt]#1\end{tcolorbox}}

\title{SPF Stellarator Neutronics: Adjoint Magnet Shielding}
\subtitle{Placement Worth Experiment}
\author{Thaddaeus Kiker}
\date{\today}

\begin{document}

\begin{frame}[plain]\titlepage\end{frame}

\begin{frame}{Since The Last Meeting}
  \begin{itemize}
    \item Last time the adjoint coil-importance map and reciprocity were validated, and the placed-versus-uniform shield test was proposed with the outcome still open.
    \item That test now resolves: nonuniform placement lowers the peak coil flux by a factor of four at equal material.
    \item I then asked the sharper question underneath it: does the cheap adjoint proxy predict where shield actually helps, patch by patch.
    \item Ground truth for twelve shield patches was measured in direct transport.
    \item Two findings came out that change how the automatic optimizer should work.
  \end{itemize}
\end{frame}

\begin{frame}{Defining Shielding Worth}
  \small
  The worth of a shield patch is the reduction in peak coil dose per unit of shield material added there, at fixed outer envelope, with breeder traded away for the shield.
  \vspace{0.4em}
  \kb{\centering $\displaystyle W_{\mathrm{FD}} \;=\; -\,\frac{\Delta(\text{coil dose})}{\Delta(\text{material})}$ \quad measured by full Monte Carlo per patch}
  \vspace{0.4em}
  \begin{itemize}
    \item Positive worth means shield placed there lowers the coil dose.
    \item The question is whether a free adjoint estimate can stand in for this expensive measurement.
  \end{itemize}
\end{frame}

\begin{frame}{Four Quantities}
  \small
  \begin{center}
  \renewcommand{\arraystretch}{1.35}
  \begin{tabular}{@{}llc@{}}
    \toprule
    Quantity & Definition & Cost \\
    \midrule
    Attribution $A$ & plasma contributon $S\,\psi^\dagger$, where neutrons are born & one adjoint \\
    Scalar worth $W$ & shield-band contributon $\phi\,\psi^\dagger$, first-order estimate & free \\
    Signed worth $W_s$ & $\sigma_{\mathrm{sh}}\!\sum_{\mathrm{shield}}\! C - \sigma_{\mathrm{br}}\!\sum_{\mathrm{breeder}}\! C$, debits breeder & free \\
    Measured worth $W_{\mathrm{FD}}$ & $-\,\Delta$dose per material, direct transport & one run per patch \\
    \bottomrule
  \end{tabular}
  \end{center}
  \vspace{0.5em}
  The first three are adjoint estimates. The fourth is the truth they are tested against.
\end{frame}

\begin{frame}{Does The Proxy Rank Placements}
  \begin{columns}[c]
    \begin{column}{0.62\textwidth}
      \includegraphics[width=\textwidth]{worth_vs_truth}
    \end{column}
    \begin{column}{0.38\textwidth}
      \footnotesize
      \begin{itemize}
        \item Across twelve patches neither the scalar worth nor the attribution orders the measured worth.
        \item Spearman correlations are $-0.34$ and $-0.17$, both consistent with no rank agreement.
        \item The cheap proxy cannot serve as the ordering an optimizer would follow.
      \end{itemize}
    \end{column}
  \end{columns}
\end{frame}

\begin{frame}{The Robust Signal Is Backfire Detection}
  \begin{columns}[c]
    \begin{column}{0.58\textwidth}
      \includegraphics[width=\textwidth]{fd_worth_bars}
    \end{column}
    \begin{column}{0.42\textwidth}
      \footnotesize
      \begin{itemize}
        \item Eleven patches reduce the coil dose by a similar amount; one raises it by forty-one percent.
        \item The scalar worth ranks that backfire patch among the best places to shield.
        \item The signed worth ranks it the worst, matching the measurement.
        \item The adjoint is reliable at the sign, not at the fine order.
      \end{itemize}
    \end{column}
  \end{columns}
\end{frame}

\begin{frame}{Why The Proxy Fails: Multigroup Moderation}
  \small
  \begin{itemize}
    \item The coil responds to fast flux. The breeder moderates fast neutrons through downscatter.
    \item Trading breeder for shield removes moderator, so fast flux is released even as shield is added.
    \item On one patch that release outweighs the added shield and the coil dose rises.
  \end{itemize}
  \vspace{0.3em}
  \kb{\centering First-order removal theory predicts dose down everywhere. The measured sign flip is a moderation effect a single total cross section cannot carry.}
  \vspace{0.3em}
  \begin{itemize}
    \item A fast-group forward flux is the correct energy weighting for the coil response.
  \end{itemize}
\end{frame}

\begin{frame}{Peak Coil Flux: Collective, Not Per-Coil}
  \begin{columns}[c]
    \begin{column}{0.56\textwidth}
      \includegraphics[width=\textwidth]{peak_flux}
    \end{column}
    \begin{column}{0.44\textwidth}
      \footnotesize
      \begin{itemize}
        \item Single-coil placement lowers the peak by a factor of four but moves it onto an unprotected coil.
        \item Distributing the same material by combined importance does not beat single-coil concentration.
        \item The coil loading is set by a shared inboard region, not by twenty independent coil footprints.
        \item Concentrating shield on that inboard region is the lever.
      \end{itemize}
    \end{column}
  \end{columns}
\end{frame}

\begin{frame}{Implications For Automatic Optimization}
  \small
  \begin{itemize}
    \item The cheap proxy does not rank placements, so it cannot replace Monte Carlo as the optimizer objective.
    \item The signed worth identifies backfire patches, so it enters as a constraint that keeps the optimizer from degrading the magnet.
    \item The collective inboard result reduces the search to a few parameters instead of a per-coil field.
    \item The workable automatic system is a low-dimensional Monte Carlo Bayesian loop with an adjoint backfire guard, not a zero-cost adjoint placer.
  \end{itemize}
\end{frame}

\begin{frame}{Reactor-Scale Coil Sets}
  \small
  \begin{itemize}
    \item Gil et al.\ 2026 released eighty-six optimized stellarator coil sets, quasi-axisymmetric and quasi-helical, from an augmented-Lagrangian method.
    \item The work is coil optimization with no neutronics, so it is a source of geometry rather than competing prior art.
    \item Coil-surface distance is an explicit parameter in the set, and that distance is our fixed blanket envelope.
    \item The sets give controlled standoff sweeps across reactor scales to test whether these results generalize.
    \item The converter into the neutronics geometry pipeline is built and tested.
  \end{itemize}
\end{frame}

\begin{frame}{In Progress}
  \begin{itemize}
    \item Concentrated placement on the shared inboard region at equal material, in transport, to test whether it beats single-coil concentration.
    \item Fast-group signed worth, to test whether the coil-relevant energy range recovers a usable ranking.
    \item Both finish tomorrow.
  \end{itemize}
\end{frame}

\end{document}
